% Figure: two panels showing the visual definition of even and odd
% functions; an even function is a mirror image across the y-axis, an
% odd function is a point reflection through the origin.
\begin{figure}[htbp]
\centering
\resizebox{\linewidth}{!}{%
\begin{tikzpicture}[scale=0.95,
  ax/.style={->, draw=engblack, thick},
  tick/.style={draw=engblack!60, thin},
  curve/.style={draw=engblue, thick},
  oddc/.style={draw=engred, thick},
  ann/.style={font=\footnotesize, align=center},
]

% --- Left panel: even, y = x^2 - 1 ---
\begin{scope}[xshift=0cm]
  \draw[ax] (-2.6, 0) -- (2.6, 0) node[right, ann] {$x$};
  \draw[ax] (0, -1.5) -- (0, 3.5) node[above, ann] {$y$};
  \foreach \x in {-2, -1, 1, 2} {
    \draw[tick] (\x, -0.05) -- (\x, 0.05);
  }
  \foreach \y in {-1, 1, 2, 3} {
    \draw[tick] (-0.05, \y) -- (0.05, \y);
  }
  \draw[curve, domain=-1.85:1.85, samples=40, smooth]
    plot (\x, {\x*\x - 1});
  \node[ann, anchor=south] at (1.65, 1.8) {$y = x^{2} - 1$};
  % Mirror cue: vertical dashed line
  \draw[dashed, draw=engblack!50] (0, -1.5) -- (0, 3.5);
  \node[ann, anchor=north, fill=white, inner sep=2pt]
    at (0, -1.7) {even: $f(-x) = f(x)$};
\end{scope}

% --- Right panel: odd, y = x^3 - 3x ---
\begin{scope}[xshift=8cm]
  \draw[ax] (-2.6, 0) -- (2.6, 0) node[right, ann] {$x$};
  \draw[ax] (0, -2.8) -- (0, 2.8) node[above, ann] {$y$};
  \foreach \x in {-2, -1, 1, 2} {
    \draw[tick] (\x, -0.05) -- (\x, 0.05);
  }
  \foreach \y in {-2, -1, 1, 2} {
    \draw[tick] (-0.05, \y) -- (0.05, \y);
  }
  \draw[oddc, domain=-2.05:2.05, samples=80, smooth]
    plot (\x, {\x*\x*\x - 3*\x});
  \node[ann, anchor=south] at (-1.5, 1.6) {$y = x^{3} - 3x$};
  % Symmetry cue: small inversion arrow through origin
  \draw[dashed, draw=engblack!50] (-2.4, -2.4) -- (2.4, 2.4);
  \node[ann, anchor=north, fill=white, inner sep=2pt]
    at (0, -3.0) {odd: $f(-x) = -f(x)$};
\end{scope}

\end{tikzpicture}%
}%
\caption[Even and odd symmetry by visual inspection]{Even functions
satisfy $f(-x) = f(x)$ and graph as a left-right mirror image across
the $y$-axis (left panel, the parabola $y = x^{2} - 1$). Odd
functions satisfy $f(-x) = -f(x)$ and graph as a point reflection
through the origin (right panel, the cubic $y = x^{3} - 3x$). An
arbitrary function decomposes uniquely into its even and odd parts,
$f(x) = \tfrac{1}{2}[f(x) + f(-x)] + \tfrac{1}{2}[f(x) - f(-x)]$, a
decomposition Volume II Chapter~13 lifts to the Fourier cosine and
sine series. Recognising parity on sight halves the integration work
for any symmetric or antisymmetric quantity (mean, kinetic energy,
overlap integral, autocorrelation).}
\label{fig:vol02:ch02:even-odd-symmetry}
\end{figure}
